
\documentclass{article}

% -----------------------------
% Packages
% -----------------------------
\usepackage{fancyhdr}
\usepackage{extramarks}
\usepackage{amsmath}
\usepackage{amsthm}
\usepackage{amsfonts}
\usepackage{tikz}
\usepackage[plain]{algorithm}
\usepackage{algpseudocode}
\usepackage{amssymb}
\usepackage{enumitem}

\usetikzlibrary{automata,positioning}

% -----------------------------
% Basic Document Settings
% -----------------------------
\topmargin=-0.45in
\evensidemargin=0in
\oddsidemargin=0in
\textwidth=6.5in
\textheight=9.0in
\headsep=0.25in

\linespread{1.1}

\pagestyle{fancy}
\lhead{\hmwkAuthorName}
\chead{\hmwkClass: \hmwkTitle}
\rhead{}
\lfoot{\lastxmark}
\cfoot{\thepage}

\renewcommand\headrulewidth{0.4pt}
\renewcommand\footrulewidth{0.4pt}

\setlength\parindent{0pt}

% -----------------------------
% Problem Section Helpers
% -----------------------------
\newcommand{\enterProblemHeader}[1]{
  \nobreak\extramarks{}{Problem #1 continued on next page\ldots}\nobreak{}
  \nobreak\extramarks{Problem #1 (continued)}{Problem #1 continued on next page\ldots}\nobreak{}
}

\newcommand{\exitProblemHeader}[1]{
  \nobreak\extramarks{Problem #1 (continued)}{Problem #1 continued on next page\ldots}\nobreak{}
  \nobreak\extramarks{Problem #1}{}\nobreak{}
}

\setcounter{secnumdepth}{0}
\newcounter{partCounter}
\newcounter{homeworkProblemCounter}
\setcounter{homeworkProblemCounter}{1}
\nobreak\extramarks{Problem \arabic{homeworkProblemCounter}}{}\nobreak{}

% -----------------------------
% Theorem env (optional)
% -----------------------------
\newtheorem*{theorem}{Theorem}

% -----------------------------
% Homework Problem Environment
%   Usage:
%   \begin{homeworkProblem}            % auto-number
%     % Your answer here
%   \end{homeworkProblem}
%
%   \begin{homeworkProblem}[1.2.3]     % label problem "1.2.3" (e.g. textbook numbering)
%     % Your answer here
%   \end{homeworkProblem}
% -----------------------------
\newenvironment{homeworkProblem}[1][]{
  \def\hwProblemLabel{#1}
  \ifx\hwProblemLabel\empty
    \edef\hwProblemLabel{\arabic{homeworkProblemCounter}}
    \stepcounter{homeworkProblemCounter}
  \fi
  \section{Problem \hwProblemLabel}
  \setcounter{partCounter}{1}
  \enterProblemHeader{\hwProblemLabel}
}{
  \exitProblemHeader{\hwProblemLabel}
}

% -----------------------------
% Homework Details (EDIT THESE)
% -----------------------------
\newcommand{\hmwkTitle}{Homework \#3}
\newcommand{\hmwkClass}{Math 3320}
\newcommand{\hmwkAuthorName}{\textbf{Joseph Quinn}}

% -----------------------------
% Title
% -----------------------------
\title{
  \vspace{2in}
  \textmd{\textbf{\hmwkClass:\ \hmwkTitle}}\\
  \vspace{3in}
}
\author{\hmwkAuthorName}
\date{}

% -----------------------------
% Convenience Macros (optional)
% -----------------------------
\renewcommand{\part}[1]{\textbf{\large Part \Alph{partCounter}}\stepcounter{partCounter}\\}

% Algorithms
\newcommand{\alg}[1]{\textsc{\bfseries \footnotesize #1}}

% Calculus
\newcommand{\deriv}[1]{\frac{\mathrm{d}}{\mathrm{d}x} (#1)}
\newcommand{\pderiv}[2]{\frac{\partial}{\partial #1} (#2)}
\newcommand{\dx}{\mathrm{d}x}

% Statistics
\newcommand{\E}{\mathrm{E}}
\newcommand{\Var}{\mathrm{Var}}
\newcommand{\Cov}{\mathrm{Cov}}
\newcommand{\Bias}{\mathrm{Bias}}

% Proof step (for aligned derivations)
\newcommand{\step}[2]{& #1 & & \text{#2} \\}

% Solution header (optional)
\newcommand{\solution}{\textbf{\large Solution}}

% -----------------------------
% Document
% -----------------------------
\begin{document}

\maketitle
\pagebreak

1.9: 5
1.11: 2,5,10 (parts a,f,h only, no need to compare to 1.11.7), 12, 13, 20.
1.12: 5,6,7,15

\begin{homeworkProblem}[1.9.5]
	$$
		\begin{array}{c|c|c|c}
			\text{Received } w & d(001,w) & d(101,w) & \text{IMLD} \\
			\hline
			000                & 1        & 2        & 001         \\
			001                & 0        & 1        & 001         \\
			010                & 2        & 3        & 001         \\
			011                & 1        & 2        & 001         \\
			100                & 2        & 1        & 101         \\
			101                & 1        & 0        & 101         \\
			110                & 3        & 2        & 101         \\
			111                & 2        & 1        & 101
		\end{array}
	$$
	Therefore IMLD correctly concludes that \(001\) was sent when $w\in\{000,001,010,011\}$.
	and  IMLD incorrectly concludes that \(101\) was sent when $w\in\{100,101,110,111\}$.
\end{homeworkProblem}

\begin{homeworkProblem}[1.11.2]
	Let

	$$
		C=\{001,101,110\}.
	$$

	\begin{enumerate}[label=\alph*.]
		\item For $u=011$,

		      $$
			      101+011=110\in C,
		      $$

		      so

		      $$
			      \boxed{C\text{ does not detect }011.}
		      $$

		\item For $u=001$,

		      $$
			      001+001=000,\qquad
			      101+001=100,\qquad
			      110+001=111.
		      $$

		      None are in $C$, so

		      $$
			      \boxed{C\text{ detects }001.}
		      $$

		\item For $u=000$,

		      $$
			      v+000=v\in C
		      $$

		      for every $v\in C$, so

		      $$
			      \boxed{C\text{ does not detect }000.}
		      $$
	\end{enumerate}
\end{homeworkProblem}

\begin{homeworkProblem}[1.11.5]

	\begin{enumerate}[label=\alph*.]
		\item If $0\in C$ and $u\in C$, choose $v=0$. Then
		      $$v+u=u\in C,$$
		      so $\boxed{C\text{ does not detect }u}$.

		\item For $u=0$,
		      $$v+0=v\in C$$
		      for every $v\in C$, so $\boxed{\text{no code detects }0}$.
	\end{enumerate}
\end{homeworkProblem}

\begin{homeworkProblem}[1.11.10]
	\begin{enumerate}[label=(\alph*)]
		\item $C=\{101,111,011\}$
		      \begin{align*}
			      101+101 & = 000 \\
			      101+111 & = 010 \\
			      101+011 & = 110 \\
			      111+011 & = 100
		      \end{align*}
		      Thus the errors that cannot be detected by $C$ are
		      \[
			      \{000,010,110,100\}.
		      \]
		      Therefore the errors that can be detected are
		      \[
			      K^3\setminus\{000,010,110,100\}.
		      \]

		\item[(f)] $C=\{00000,11100,00111,11011\}$
		      \begin{align*}
			      00000+00000 & = 00000 \\
			      00000+11100 & = 11100 \\
			      00000+00111 & = 00111 \\
			      00000+11011 & = 11011 \\
			      11100+00111 & = 11011 \\
			      11100+11011 & = 00111 \\
			      00111+11011 & = 11100
		      \end{align*}
		      Thus the errors that cannot be detected by $C$ are
		      \[
			      \{00000,11100,00111,11011\}.
		      \]
		      Therefore the errors that can be detected are
		      \[
			      K^5\setminus\{00000,11100,00111,11011\}.
		      \]
		\item[(h)]
		      $C=\{000000,101010,010101,111111\}$
		      \begin{align*}
			      000000+000000 & = 000000 \\
			      000000+101010 & = 101010 \\
			      000000+010101 & = 010101 \\
			      000000+111111 & = 111111 \\
			      101010+101010 & = 000000 \\
			      101010+010101 & = 111111 \\
			      101010+111111 & = 010101 \\
			      010101+010101 & = 000000 \\
			      010101+111111 & = 101010 \\
			      111111+111111 & = 000000
		      \end{align*}
		      Thus the errors that cannot be detected by $C$ are
		      \[
			      \{000000,101010,010101,111111\}.
		      \]
		      Therefore the errors that can be detected are
		      \[
			      K^6\setminus\{000000,101010,010101,111111\}.
		      \]
	\end{enumerate}
\end{homeworkProblem}

\begin{homeworkProblem}[1.11.12]
	\begin{enumerate}[label=(\alph*)]
		\item $C=\{101,111,011\}$

		      \[
			      d(101,111)=1,\quad d(101,011)=2,\quad d(111,011)=1,
		      \]
		      so $d(C)=1$.

		\item $C=\{000,001,010,011\}$

		      \[
			      d(000,001)=1,
		      \]
		      so $d(C)=1$.

		\item $C=\{0000,0001,1110\}$

		      \[
			      d(0000,0001)=1,
		      \]
		      so $d(C)=1$.

		\item $C=\{0000,1001,0110,1111\}$

		      \[
			      d(0000,1001)=2,
		      \]
		      so $d(C)=2$.

		\item $C=\{00000,11111\}$

		      \[
			      d(00000,11111)=5,
		      \]
		      so $d(C)=5$.

		\item $C=\{00000,11100,00111,11011\}$

		      \[
			      d(00000,11100)=3,
		      \]
		      so $d(C)=3$.

		\item $C=\{00000,11110,01111,10001\}$

		      \[
			      d(00000,10001)=2,
		      \]
		      so $d(C)=2$.

		\item $C=\{000000,101010,010101,111111\}$

		      \[
			      d(000000,101010)=3,
		      \]
		      so $d(C)=3$.
	\end{enumerate}
\end{homeworkProblem}
\begin{homeworkProblem}[1.11.13]
	Adding a parity check digit to $K^n$ gives a code in which every codeword has even weight.

	If two words in $K^n$ differ in exactly one position, then their parity check digits must also differ. Thus the corresponding codewords differ in two positions.

	Therefore the minimum distance is

	\[
		d(C)=2.
	\]
	\\
	example with $K^2 = \{00,01,10,11\}$
	after adding parity check, $C = \{000,011,101,110\} $ and $d(C) = 2$
\end{homeworkProblem}
\begin{homeworkProblem}[1.11.20]
	If we think of even and odd numbers we know an even number plus an odd number will never equal an even number. \\
	Taking that same logic if we et $C$ be the code consisting of all words of length $4$ with even weight.
	If $u$ has odd weight, then for every $v\in C$, the word
	\[
		v+u
	\]
	has odd weight, so $v+u\notin C$. Thus $C$ detects every odd weight error pattern.

	If $u$ has even weight, then for every $v\in C$, the word $v+u$ also has even weight, so $v+u\in C$. Thus $C$ does not detect any even weight error pattern.

	Therefore $C$ detects exactly the odd weight error patterns in $k^4$:
	\[
		\{0001,0010,0100,1000,1110,1101,1011,0111\}.
	\]
\end{homeworkProblem}

\begin{homeworkProblem}[1.12.5]
	Let
	\[
		C=\{001,101,110\}.
	\]

	\begin{enumerate}[label=\alph*)]
		\item $u=100$

		      Take $v=001$. Then
		      \[
			      w=v+u=001+100=101.
		      \]

		      Compare $w$ to each codeword:
		      \begin{align*}
			      d(001,w) & = d(001,101)=1 \\
			      d(101,w) & = d(101,101)=0 \\
			      d(110,w) & = d(110,101)=2
		      \end{align*}

		      Thus IMLD chooses $101$ instead of the transmitted word $001$.
		      Therefore,
		      \[
			      \boxed{C \text{ does not correct }100.}
		      \]

		\item $u=000$

		      \begin{itemize}
			      \item $v=001$, so $w=001+000=001$
			            \begin{align*}
				            d(001,w) & = d(001,001)=0 \\
				            d(101,w) & = d(101,001)=1 \\
				            d(110,w) & = d(110,001)=3
			            \end{align*}
			            Thus IMLD chooses $001$.

			      \item $v=101$, so $w=101+000=101$
			            \begin{align*}
				            d(001,w) & = d(001,101)=1 \\
				            d(101,w) & = d(101,101)=0 \\
				            d(110,w) & = d(110,101)=2
			            \end{align*}
			            Thus IMLD chooses $101$.

			      \item $v=110$, so $w=110+000=110$
			            \begin{align*}
				            d(001,w) & = d(001,110)=3 \\
				            d(101,w) & = d(101,110)=2 \\
				            d(110,w) & = d(110,110)=0
			            \end{align*}
			            Thus IMLD chooses $110$.
		      \end{itemize}

		      Therefore,
		      \[
			      \boxed{C \text{ corrects }000.}
		      \]
	\end{enumerate}
\end{homeworkProblem}

\begin{homeworkProblem}[1.12.6]
	\begin{proof}
		suppose the same error pattern appeared in a single row twice, meaning we had
		two valid codewords $v_{1}$ and $v_{2}$ where $v_{1}+u=w$ and $v_{2}+u=w$ \\
		\begin{align*}
			v_{1}+u=v_{2}+u     \\
			v_{1}+u+u=v_{2}+u+u \\
			v_{1}+0 = v_{2}+0   \\
			\text{so} \; v_{1} = v_{2}
		\end{align*}
		Therefore the same error pattern cannot occur more than once in a given row.
	\end{proof}

\end{homeworkProblem}

\begin{homeworkProblem}[1.12.7]
	\begin{proof}
		Let's take any $v \in C$.

		\[
			d(v,v+0)=0
		\]

		For any other codeword $w\in C$ where $w\neq v$,

		\[
			d(w,v+0)=d(w,v)>0.
		\]

		Thus,

		\[
			d(v,v+0)<d(w,v+0)
			\qquad \forall w\in C,\; w\neq v.
		\]

		Therefore $C$ corrects the error pattern $0$.
	\end{proof}
\end{homeworkProblem}

\begin{homeworkProblem}[1.12.15]
	Let
	\[
		C=\{0000,0011,0101,0110,1001,1010,1100,1111\}.
	\]

	The minimum distance of $C$ is
	\[
		d(C)=2.
	\]

	From the theorem we know $C$ corrects all error patterns of weight at most
	\[
		\left\lfloor \frac{d-1}{2} \right\rfloor
		=
		\left\lfloor \frac{1}{2} \right\rfloor
		=0.
	\]

	Meaning $C$ will only correct $0000$

	But lets say we let $u\neq 0000$ two things could happen,
	\begin{enumerate}
		\item
		      If $u$ has even weight, then for any $v\in C$, $v+u$ also has even weight, so
		      \[
			      v+u\in C.
		      \]
		      Since $u\neq 0$, we have $v+u\neq v$, so $u$ is not corrected.

		\item

		      If $u$ has odd weight, then $v+u$ has odd weight. Any odd-weight word of length $4$ is distance $1$ from more than one even-weight codeword, so IMLD cannot uniquely decode it back to $v$.

		      Therefore the only error pattern corrected by $C$ is $0000$
	\end{enumerate}
\end{homeworkProblem}
\end{document}
